// Film Chromatic Aberration — multi-band spectral lateral CA
//
// Simulates real lateral chromatic aberration by decomposing the image
// into N spectral bands, applying wavelength-dependent radial displacement
// to each, and reintegrating via spectral-to-RGB color matching weights.
//
// This produces the smooth rainbow fringing of real glass — unlike the
// crude "shift R and B channels" trick which creates harsh cyan/red banding.
//
// Model: each spectral band lambda_i receives a radial scale:
//   r'(lambda) = r * (1 + alpha * lambda_i * r^p)
// where lambda_i ranges from -1 (blue) to +1 (red), alpha is the global
// CA strength, and p is the radial falloff exponent (2 = quadratic).
// Green (lambda=0) is held fixed as the reference channel.
//
// The spectral bands are weighted by approximate CIE color matching
// functions to reconstruct physically plausible RGB.
//
// Connect @image in scene-linear color space.

f$strength = 0.003;     // lateral CA intensity (start small: 0.001-0.01)
f$falloff = 2.0;        // radial exponent (2 = quadratic, 3 = cubic)
i$bands = 8;            // spectral samples (8-16 for smooth fringing)
f$center_x = 0.5;       // optical center X
f$center_y = 0.5;       // optical center Y

// ── Setup ────────────────────────────────────────────────────────
float aspect = iw / ih;
float cx = (u - $center_x) * aspect;
float cy = v - $center_y;
float r = hypot(cx, cy);
float r_pow = pow(max(r, 0.0001), $falloff);

// Direction from center (unit vector for radial displacement)
float inv_r = sdiv(1.0, r);
float dir_x = cx * inv_r;
float dir_y = cy * inv_r;

// Clamp bands: minimum 3 for smooth spectrum, maximum 64 (well within loop limit)
int num_bands = clamp($bands, 3, 64);
float inv_bands = 1.0 / float(num_bands - 1);

// ── Accumulate spectral bands ────────────────────────────────────
// Each band gets a wavelength parameter lambda in [-1, +1].
// Approximate CIE color matching: R peaks at lambda=+0.7, G at 0,
// B at lambda=-0.7, with overlapping Gaussian response curves.
vec3 accum = vec3(0.0);
vec3 weight_sum = vec3(0.0);

for (int i = 0; i < num_bands; i++) {
    // Normalized wavelength: -1 (blue end) to +1 (red end)
    float lambda = float(i) * inv_bands * 2.0 - 1.0;

    // Radial scale for this wavelength
    float displacement = $strength * lambda * r_pow;
    float su = (cx + dir_x * displacement) / aspect + $center_x;
    float sv = (cy + dir_y * displacement) + $center_y;

    vec3 s = sample(@image, clamp(su, 0.0, 1.0), clamp(sv, 0.0, 1.0));

    // Spectral-to-RGB weights (Gaussian approximation of CIE curves)
    // Red:   peak at lambda = +0.65, sigma = 0.45
    // Green: peak at lambda =  0.00, sigma = 0.40
    // Blue:  peak at lambda = -0.65, sigma = 0.40
    float wr = exp(-0.5 * pow((lambda - 0.65) / 0.45, 2.0));
    float wg = exp(-0.5 * pow(lambda / 0.40, 2.0));
    float wb = exp(-0.5 * pow((lambda + 0.65) / 0.40, 2.0));

    accum = accum + vec3(s.r * wr, s.g * wg, s.b * wb);
    weight_sum = weight_sum + vec3(wr, wg, wb);
}

// Energy-preserving normalization
vec3 result = vec3(
    sdiv(accum.r, weight_sum.r),
    sdiv(accum.g, weight_sum.g),
    sdiv(accum.b, weight_sum.b)
);

@OUT = result;
